% Luc Destombe --- licence CC BY-NC-SA 4.0
% https://creativecommons.org/licenses/by-nc-sa/4.0/deed.fr
% Fichier autonome : compiler avec pdflatex, deux fois (signets, nombre de pages).
\documentclass[11pt,a4paper]{article}
\makeatletter
\newif\ifeleve
\elevefalse

\newif\iflycee@palatino
\lycee@palatinotrue

\RequirePackage[utf8]{inputenc}
\RequirePackage[T1]{fontenc}
\RequirePackage[french]{babel}
\RequirePackage{etoolbox}

\newcommand{\ShorthandsOff}{\@ifundefined{shorthandoff}{}{\shorthandoff{;:!?}}}
\newcommand{\ShorthandsOn}{\@ifundefined{shorthandon}{}{\shorthandon{;:!?}}}
\ShorthandsOff
\AtBeginDocument{\ShorthandsOn}
\AtBeginEnvironment{tikzpicture}{\ShorthandsOff}

\RequirePackage{geometry}
\RequirePackage{fancyhdr}
\RequirePackage{lastpage}
\RequirePackage{multicol}
\RequirePackage{array}
\RequirePackage{enumitem}
\RequirePackage{xcolor}
\RequirePackage{needspace}

\RequirePackage{amsmath,amssymb}
\RequirePackage{mathpazo}
\DeclareMathSizes{10.95}{10.95}{8}{5.5}
\begingroup\catcode`\_=\active \catcode`\^=\active
\gdef_#1{\sb{\scriptscriptstyle #1}}
\gdef^#1{\sp{\scriptscriptstyle\lycee@moinsepais #1}}
\endgroup
\newcommand\lycee@moins{\mathbin{%
  \setbox\z@\hbox{$\m@th\scriptscriptstyle\lycee@moinsorig$}%
  \dimen@=.55\wd\z@ \advance\dimen@-.5pt
  \hbox to.75\wd\z@{\hss$\m@th\scriptscriptstyle\vcenter{\hbox{%
    \tikz\draw[line width=.5pt,line cap=round](0,0)--(\the\dimen@,0);}}$\hss}}}
\begingroup\lccode`\~=`\- \lowercase{\endgroup\let~}\lycee@moins
\newcommand\lycee@moinsepais{\mathcode`\-="8000 }
\AtBeginDocument{\mathchardef\lycee@moinsorig=\mathcode`\-\relax}
\AtBeginDocument{\catcode`\_=12 \mathcode`\_="8000
  \catcode`\^=12 \mathcode`\^="8000 }
\newcommand\lycee@exposants{%
  \fontdimen13\textfont2=.48\fontdimen6\textfont2
  \fontdimen14\textfont2=.48\fontdimen6\textfont2
  \fontdimen15\textfont2=.48\fontdimen6\textfont2
  \fontdimen13\scriptfont2=.48\fontdimen6\scriptfont2
  \fontdimen14\scriptfont2=.48\fontdimen6\scriptfont2
  \fontdimen15\scriptfont2=.48\fontdimen6\scriptfont2}
\every@math@size\expandafter{\the\every@math@size\lycee@exposants}
\newlength{\retraitcalculs}
\setlength{\retraitcalculs}{2em}

\RequirePackage{tikz}
\usetikzlibrary{arrows.meta,calc,babel,shapes.misc}
\RequirePackage[most]{tcolorbox}
\tcbuselibrary{skins,breakable}

\definecolor{coulDef}{HTML}{1F4E79}
\definecolor{coulProp}{HTML}{7030A0}
\definecolor{coulRem}{HTML}{C55A11}
\definecolor{coulEx}{HTML}{2E7D32}
\definecolor{coulExo}{HTML}{B03A2E}
\definecolor{coulPart}{HTML}{1F4E79}
\definecolor{coulMeth}{HTML}{1B6E4F}
\definecolor{coulCourbe}{HTML}{0B5ED7}
\definecolor{coulCourbeB}{HTML}{C00000}
\definecolor{coulCourbeC}{HTML}{2E7D32}
\definecolor{coulGrille}{HTML}{8C8C8C}
\definecolor{coulRep}{HTML}{1B6E4F}
\definecolor{coulBar}{HTML}{2E7D32}

\newcommand{\R}{\mathbb{R}}
\newcommand{\Rs}{\mathbb{R}^{*}}
\renewcommand{\le}{\leqslant}
\renewcommand{\ge}{\geqslant}
\newcommand{\vect}[1]{\overrightarrow{#1}}
\newcommand{\norme}[1]{\left\lVert #1 \right\rVert}
\newcommand{\intervalle}[2]{\left[#1\,;#2\right]}
\RequirePackage{cancel}
\renewcommand{\CancelColor}{\color{coulExo}}
\newcommand{\fois}[1]{\times{\color{coulExo}#1}}
\newcommand{\barre}[1]{\cancel{#1}}

\tikzset{grille/.style={coulGrille,line width=0.4pt}}
\newcommand{\quadrillage}[4]{\draw[grille,step=1] (#1,#2) grid (#3,#4);}
\tikzset{
  croix/.style={cross out,draw,line width=0.6pt,inner sep=0pt,minimum size=4pt},
  croixdroite/.style={croix,rotate=45,line width=0.8pt},
}
\newcommand{\pt}[3]{\path (#1) node[croix]{} node[#2,font=\small,inner sep=2.5pt]{$#3$};}

\newcommand{\lycee@repere}[4]{%
  \draw[grille,step=1] (#1,#2) grid (#3,#4);
  \draw[-{Stealth[length=2mm]},thick] (#1,0)--({#3+0.4},0) node[below left=-1pt]{$x$};
  \draw[-{Stealth[length=2mm]},thick] (0,#2)--(0,{#4+0.4}) node[below left=-1pt]{$y$};
  \node[below left,font=\scriptsize,inner sep=1.5pt] at (0,0) {$0$};}
\newcommand{\lycee@gradx}[1]{\foreach \k/\l in {#1}{%
  \draw (\k,0.07)--(\k,-0.07) node[below,font=\scriptsize,inner sep=1.5pt]{$\l$};}}
\newcommand{\lycee@grady}[1]{\foreach \k/\l in {#1}{%
  \draw (0.07,\k)--(-0.07,\k) node[left,font=\scriptsize,inner sep=1.5pt]{$\l$};}}
\newcommand{\lycee@intff}[2]{\left[#1\,;#2\right]}
\newcommand{\lycee@intfo}[2]{\left[#1\,;#2\right[}
\newcommand{\lycee@intof}[2]{\left]#1\,;#2\right]}
\newcommand{\lycee@intoo}[2]{\left]#1\,;#2\right[}
\newcommand{\lycee@ens}[1]{\left\{#1\right\}}
\AtBeginDocument{%
  \providecommand{\repere}{\lycee@repere}%
  \providecommand{\gradx}{\lycee@gradx}%
  \providecommand{\grady}{\lycee@grady}%
  \providecommand{\Cf}{\mathcal{C}\sb{\scriptscriptstyle f}}%
  \providecommand{\Cg}{\mathcal{C}\sb{\scriptscriptstyle g}}%
  \providecommand{\intff}{\lycee@intff}%
  \providecommand{\intfo}{\lycee@intfo}%
  \providecommand{\intof}{\lycee@intof}%
  \providecommand{\intoo}{\lycee@intoo}%
  \providecommand{\ens}{\lycee@ens}}

\newlist{questions}{enumerate}{3}
\setlist[questions,1]{label=\arabic*\textdegree),leftmargin=*,itemsep=2pt,topsep=3pt}
\setlist[questions,2]{label=\alph*),leftmargin=*,itemsep=1pt,topsep=2pt}
\setlist[questions,3]{label=\roman*),leftmargin=*,itemsep=1pt,topsep=1pt}
\setlist[itemize]{label=\textbullet,leftmargin=*,itemsep=2pt,topsep=3pt}

\newcommand{\besoinplace}[1]{\par\penalty\@M\begingroup
  \dimen@=#1\relax\dimen@ii=\pagegoal\advance\dimen@ii-\pagetotal
  \ifdim\dimen@>\dimen@ii\ifdim\dimen@ii>\z@\vfil\fi\break\fi\endgroup}

\newcommand{\lycee@pied}{\thepage/\pageref{LastPage}}
\newcommand{\entete}[2]{%
  \pagestyle{fancy}\fancyhf{}%
  \fancyhead[L]{\footnotesize\color{coulPart}\textbf{#1}}%
  \fancyhead[R]{\footnotesize\color{coulPart}\textbf{#2}}%
  \fancyfoot[C]{\footnotesize\lycee@pied}%
  \renewcommand{\headrulewidth}{0.4pt}%
  \renewcommand{\headrule}{\hbox to\headwidth{%
    \color{coulPart!40}\leaders\hrule height \headrulewidth\hfill}}}

\RequirePackage{ccicons}
\newcommand{\lycee@licence}{{\scriptsize\color{gray}%
  \copyright\ Luc Destombe --- \ccbyncsa\ CC BY-NC-SA 4.0}}
\newif\iflycee@licence \lycee@licencetrue
\newcommand{\sanslicence}{\lycee@licencefalse}
\AtBeginDocument{\iflycee@licence\AddToHookNext{shipout/foreground}{%
  \put(0,0){\rlap{\kern\dimexpr1in+\hoffset+\oddsidemargin+\textwidth\relax
    \llap{\raisebox{-\dimexpr1in+\voffset+\topmargin+\headheight+\headsep
      +\textheight+\footskip\relax}{\lycee@licence}}}}}\fi}

\newcommand{\formulesserrees}{%
  \setlength{\abovedisplayskip}{3pt}\setlength{\belowdisplayskip}{3pt}%
  \setlength{\abovedisplayshortskip}{2pt}\setlength{\belowdisplayshortskip}{3pt}}

\setlength{\parindent}{0pt}

\RequirePackage[implicit=false,hidelinks,pdfencoding=auto,psdextra,bookmarksopen,
  bookmarksopenlevel=1,pdfcreator={},pdfproducer={}]{hyperref}
\RequirePackage{bookmark}
\hypersetup{pdfauthor={Luc Destombe},pdfkeywords={CC BY-NC-SA 4.0}}
\newcounter{lycee@signet}
\newcommand{\signet}[3][]{\stepcounter{lycee@signet}%
  \edef\lycee@dest{\if\relax\detokenize{#1}\relax
    signet.\arabic{lycee@signet}\else#1\fi}%
  \hypertarget{\lycee@dest}{}\bookmark[level=#2,dest=\lycee@dest]{#3}}
\pdfstringdefDisableCommands{%
  \def\R{R}\def\vect#1{#1}\def\Cf{Cf}\def\Cg{Cg}\def\fois#1{×#1}%
  \def\barre#1{#1}\def\intervalle#1#2{[#1;#2]}\def\intff#1#2{[#1;#2]}%
  \def\textsuperscript#1{#1}\def\dfrac#1#2{#1/#2}\def\frac#1#2{#1/#2}%
  \def\vec#1{#1⃗}}%
\geometry{top=2cm,bottom=1cm,left=1cm,right=1cm,headheight=14pt,footskip=0.6cm}

\RequirePackage{pgfplots}
\pgfplotsset{compat=1.18}
\RequirePackage{tkz-tab}

\newcounter{partiecours}
\renewcommand{\thepartiecours}{\Roman{partiecours}}
\newcommand{\partiecours}[1]{%
  \refstepcounter{partiecours}%
  \Needspace*{8\baselineskip}%
  \bigskip
  \noindent\begin{tcolorbox}[enhanced,colback=coulPart,colframe=coulPart,
      boxrule=0pt,arc=2pt,left=8pt,right=8pt,top=4pt,bottom=4pt,
      before skip=6pt,after skip=10pt]
    \color{white}\bfseries\large \thepartiecours{} --- #1\signet{1}{\thepartiecours{} --- #1}
  \end{tcolorbox}%
  \addcontentsline{toc}{section}{\thepartiecours{} --- #1}%
}

\newtcolorbox{boitecours}[3][]{%
  enhanced,breakable,
  colback=white,colframe=#2,coltitle=white,
  fonttitle=\bfseries,
  attach boxed title to top left={xshift=8pt,yshift=-\tcboxedtitleheight/2},
  boxed title style={colback=#2,arc=2pt,boxrule=0pt},
  boxrule=0.9pt,arc=2pt,
  left=8pt,right=8pt,top=8pt,bottom=6pt,
  title={#3},#1}

\newcounter{methode}
\newenvironment{definitionbox}[1][Définition]%
  {\begin{boitecours}[phantom={\signet{2}{#1}}]{coulDef}{#1}}{\end{boitecours}}
\newenvironment{proprietebox}[1][Propriété]%
  {\begin{boitecours}[phantom={\signet{2}{#1}}]{coulProp}{#1}}{\end{boitecours}}
\newenvironment{methodebox}[1][Méthode]{\stepcounter{methode}%
  \begin{boitecours}[phantom={\signet[methode-\arabic{methode}]{2}{#1}}]{coulMeth}{#1}}%
  {\end{boitecours}}

\newcommand{\etape}[1]{\par\smallskip\textbf{\color{coulMeth}#1}\par\nobreak\smallskip}
\newcommand{\enonce}[1]{\emph{#1}\par\smallskip}
\newcommand{\reponse}[1]{\par\smallskip
  {\footnotesize\itshape\color{black!55}Réponses : #1}}
\newcommand{\demo}[1]{\textbf{\color{coulMeth}Démonstration.} #1}
\newcommand{\afaire}[1]{%
  \par\smallskip
  \noindent{\small\color{coulExo}$\blacktriangleright$\ \textbf{À faire :}
    fiche d'exercices du chapitre, #1.}\par\smallskip}

\newenvironment{calculs}{\setlength{\jot}{5pt}\@fleqntrue\@mathmargin\retraitcalculs
  \setlength{\abovedisplayskip}{4pt}\setlength{\belowdisplayskip}{4pt}%
  \setlength{\abovedisplayshortskip}{4pt}\setlength{\belowdisplayshortskip}{4pt}%
  \csname alignat*\endcsname{2}}%
  {\csname endalignat*\endcsname}
\newcommand{\rem}[1]{\text{\small\itshape\color{black!60}#1}}
\newcommand{\explique}[2]{\par\smallskip\noindent
  \begin{minipage}[t]{0.57\linewidth}\raggedright #1\end{minipage}\hfill
  \begin{minipage}[t]{0.40\linewidth}\raggedright\leavevmode
    {\small\itshape\color{black!60}#2}\end{minipage}\par}
\newcommand{\cas}[1]{\par\smallskip\noindent\textbf{\color{coulExo}#1}\nobreak}
\newcommand{\calc}[1]{\par\smallskip\textbf{\color{coulExo}#1}\par\nobreak}

\newtcolorbox{remarquebox}[1][Remarque]{%
  enhanced,breakable,colback=white,colframe=coulRem,
  boxrule=0pt,leftrule=2.5pt,arc=0pt,outer arc=0pt,
  left=8pt,right=6pt,top=5pt,bottom=5pt,
  before upper={\textbf{\color{coulRem}#1}\par\smallskip}}

\newtcolorbox{exemplebox}[1][Exemple]{%
  enhanced,breakable,colback=white,colframe=coulEx,
  boxrule=0pt,leftrule=2.5pt,arc=0pt,outer arc=0pt,
  left=8pt,right=6pt,top=5pt,bottom=5pt,
  before upper={\textbf{\color{coulEx}#1}\par\smallskip}}

\pgfplotsset{
  repere/.style={
    axis lines=middle,
    axis line style={-{Stealth[length=2.4mm]},thick},
    every axis x label/.style={at={(ticklabel* cs:1.0)},anchor=west},
    every axis y label/.style={at={(ticklabel* cs:1.0)},anchor=south},
    label style={font=\footnotesize},
    tick label style={font=\scriptsize},
    grid=both,
    major grid style={grille},
    minor tick num=0,
    tick align=inside,
    clip mode=individual,
  },
}
\tikzset{
  courbe/.style={coulCourbe,line width=1.1pt,smooth},
  courbeB/.style={coulCourbeB,line width=1.1pt,smooth},
  courbeC/.style={coulCourbeC,line width=1.1pt,smooth},
}

\newlist{etapes}{enumerate}{1}
\setlist[etapes]{label=\textbf{\arabic*.},leftmargin=*,itemsep=1pt,topsep=2pt}

\newlist{expressions}{enumerate}{1}
\setlist[expressions]{label={\Alph*\ $=$},leftmargin=*,itemsep=3pt,topsep=3pt}

\newcounter{exo}
\renewcommand{\theexo}{\ifnum\value{exo}<10 0\fi\arabic{exo}}
\newtcolorbox{exercice}[1][]{%
  enhanced,breakable,
  colback=white,colframe=coulExo,
  boxrule=0.9pt,arc=2pt,
  left=8pt,right=8pt,top=8pt,bottom=6pt,
  coltitle=white,fonttitle=\bfseries,
  attach boxed title to top left={xshift=8pt,yshift=-\tcboxedtitleheight/2},
  boxed title style={colback=coulExo,arc=2pt,boxrule=0pt},
  phantom={\signet{2}{Exercice \arabic{exo}}},
  title={Exercice \theexo\ifx\relax#1\relax\else{} --- #1\fi}}
\newenvironment{exo}[1][\relax]{\refstepcounter{exo}\begin{exercice}[#1]}{\end{exercice}}

  \newenvironment{correction}{\par}{\par}
  \newcommand{\autableau}{}
  \newcommand{\etapeexemples}{\etape{Exemples corrigés}}
  \newcommand{\etapetableau}[1]{\etape{#1}}
  \newenvironment{exempletableau}[1][Exemple]%
    {\begin{exemplebox}[#1]}{\end{exemplebox}}

\newcommand{\coursEntete}{}
\newcommand{\coursNiveau}{}
\newcommand{\coursTitre}{}
\newcommand{\coursSurTitre}{}
\newcommand{\coursSousTitre}{}

\newcommand{\enteteCours}[1]{\renewcommand{\coursEntete}{#1}}
\newcommand{\chapitrecours}[2]{\enteteCours{Chapitre #1 --- #2}}
\newcommand{\niveaucours}[1]{\renewcommand{\coursNiveau}{#1}}
\newcommand{\titrecours}[3][]{%
  \renewcommand{\coursTitre}{#2}%
  \renewcommand{\coursSurTitre}{#1}%
  \renewcommand{\coursSousTitre}{#3}}

\pagestyle{fancy}
\fancyhf{}
\fancyhead[L]{\footnotesize\color{coulPart}\textbf{\coursEntete}}
\fancyhead[R]{\footnotesize\color{coulPart}\textbf{\coursNiveau}}
\fancyfoot[C]{\footnotesize\thepage}
\renewcommand{\headrulewidth}{0.4pt}
\renewcommand{\headrule}{\hbox to\headwidth{\color{coulPart!40}\leaders\hrule height \headrulewidth\hfill}}

\setlength{\parskip}{2pt}

\AtBeginDocument{%
  \begin{center}
    \begin{tcolorbox}[enhanced,width=0.72\linewidth,colback=white,colframe=coulPart,
        boxrule=1.6pt,arc=3pt,halign=center,top=10pt,bottom=10pt]
      \ifx\coursSurTitre\empty
        {\Huge\bfseries\color{coulPart} \coursTitre}\\[4pt]
      \else
        {\Huge\bfseries\color{coulPart} \coursTitre}\\[3pt]
        {\Large\bfseries\color{coulPart} \coursSurTitre}\\[5pt]
      \fi
      {\normalsize \coursSousTitre}%
      
    \end{tcolorbox}
  \end{center}%
}
\makeatother

\chapitrecours{2}{Fonctions}
\niveaucours{1\textsuperscript{re} STI2D}
\titrecours{FONCTIONS}{Cours --- Première STI2D}

\definecolor{coulCourbe}{HTML}{C0392B}
\definecolor{coulCourbeB}{HTML}{1F4E79}

\begin{document}

\begin{remarquebox}[Comment utiliser ce cours]
Chaque partie commence par ce qu'il faut \textbf{savoir} (définitions, théorèmes), puis
vient ce qu'il faut \textbf{savoir faire} : une méthode, avec des
\textbf{exemples corrigés} faits en classe, puis des exercices
\og\textbf{À vous de jouer}\fg{} à chercher seul.

\end{remarquebox}

\partiecours{Antécédent et image}

\begin{definitionbox}[Définition 1 --- Fonction, image, antécédent]
Une \textbf{fonction} $f$ associe à chaque nombre $x$ \textbf{un seul} nombre, noté
$f(x)$.
\begin{center}
\begin{tikzpicture}[>={Stealth[length=3mm]}]
  \node (x) at (0,0) {$x$};
  \node (y) at (4,0) {$y=f(x)$};
  \draw[->,line width=1pt,coulEx] (x) to[bend left=30]
    node[above,font=\scriptsize\bfseries,coulEx] {$f$} (y);
\end{tikzpicture}
\end{center}
\begin{itemize}
  \item $f(x)$ est l'\textbf{image} de $x$ : il y en a \textbf{une seule} ;
  \item $x$ est \textbf{un antécédent} de $f(x)$ : un nombre peut avoir
        \textbf{plusieurs} antécédents, ou \textbf{aucun}.
\end{itemize}
On écrit aussi $f:x\mapsto x^{2}-4x$ pour $f(x)=x^{2}-4x$. $f$ est la fonction, $f(x)$
est un nombre.
\end{definitionbox}

\begin{methodebox}[Méthode 1 --- Calculer une image, tester un antécédent]
\etapeexemples
On considère la fonction $f$ définie par $f(x)=x^{2}-4x$.
\begin{questions}[label=\alph*)]
  \item Calculer l'image de $2$, puis celle de $-1$.
  \item Le nombre $5$ est-il un antécédent de $5$ par $f$ ?
\end{questions}
\begin{correction}
\cas{a)}
\begin{calculs}
f(2) &= 2^{2}-4\times 2 &\qquad& \rem{on remplace $x$ par $2$}\\
f(2) &= -4 &&\\[2pt]
f(-1) &= (-1)^{2}-4\times(-1) &\qquad& \rem{un nombre négatif se met entre parenthèses}\\
f(-1) &= 1+4 &\qquad& \rem{$(-1)^{2}=1$, alors que $-1^{2}=-1$}\\
f(-1) &= 5 &&
\end{calculs}
\cas{b)}
\begin{calculs}
f(5) &= 5^{2}-4\times 5 &\qquad& \rem{« $5$ est un antécédent de $5$ » veut dire $f(5)=5$ : on calcule}\\
f(5) &= 5 &&
\end{calculs}
\explique{Oui : $f(5)=5$, donc $5$ est un antécédent de $5$.}{d'après a), $-1$ en est un
autre : un nombre peut avoir plusieurs antécédents}
\end{correction}

\etape{À vous de jouer}
On considère la fonction $g$ définie par $g(x)=4x-9$.
\begin{questions}
  \item Calculer l'image de $-2$, de $5$, puis de $\dfrac{3}{4}$.
  \item Le nombre $\dfrac{1}{2}$ est-il un antécédent de $-7$ par $g$ ?
\end{questions}
\reponse{1\textdegree) $g(-2)=-17$ ; $g(5)=11$ ; $g\!\left(\frac{3}{4}\right)=-6$ ;\;
2\textdegree) $g\!\left(\frac{1}{2}\right)=2-9=-7$ : oui.}
\end{methodebox}

\afaire{exercices 1 à 4}

\partiecours{Ensemble de définition}

\begin{definitionbox}[Définition 2 --- Ensemble de définition]
L'\textbf{ensemble de définition} de $f$, noté $D_{f}$, est l'ensemble des nombres $x$
pour lesquels on peut calculer $f(x)$.
\end{definitionbox}

\begin{proprietebox}[Propriété 1 --- Deux calculs impossibles]
\begin{itemize}
  \item On ne divise pas par $0$ : un \textbf{dénominateur} ne doit pas être nul.
  \item La \textbf{racine carrée} d'un nombre négatif n'existe pas : sous une racine,
        il faut un nombre positif ou nul.
\end{itemize}
Sans fraction ni racine, on peut toujours calculer $f(x)$ : $D_{f}=\R$.
\end{proprietebox}

\begin{methodebox}[Méthode 2 --- Déterminer un ensemble de définition]
\etapeexemples
Déterminer l'ensemble de définition de chaque fonction.
\begin{questions}[label=\alph*)]
  \item $f(x)=\dfrac{x-1}{x+4}$
  \item $g(x)=\dfrac{7}{3x-6}$
  \item $h(x)=\sqrt{x+3}$
  \item $k(x)=\sqrt{2x-8}$
\end{questions}
\begin{correction}
\cas{a)}
\explique{$x+4=0$ pour $x=-4$ : $D_{f}=\R\setminus\{-4\}$.}{on cherche la valeur qui
annule le dénominateur ; « $\R$ privé de $-4$ »}
\cas{b)}
\begin{calculs}
3x-6 &= 0 &&\\
3x &= 6 &\qquad& \rem{valeur qui annule le dénominateur}\\
       x &= 2 &&
\end{calculs}
\explique{$D_{g}=\R\setminus\{2\}$.}{}
\cas{c)}
\begin{calculs}
x+3 &\ge 0 &&\\
x &\ge -3 &\qquad& \rem{sous la racine : positif ou nul}
\end{calculs}
\explique{$D_{h}=\left[-3\,;+\infty\right[$.}{$-3$ compris : $\sqrt{0}=0$ existe}
\cas{d)}
\begin{calculs}
2x-8 &\ge 0 &&\\
2x &\ge 8 &&\\
x &\ge 4 &&
\end{calculs}
\explique{$D_{k}=\left[4\,;+\infty\right[$.}{contrôle : $k(6)=\sqrt{4}=2$ existe,
$k(0)=\sqrt{-8}$ n'existe pas}
\end{correction}

\etape{À vous de jouer}
Déterminer l'ensemble de définition de chaque fonction :
\[
  f(x)=\dfrac{5}{x-6} \qquad g(x)=\dfrac{x+2}{5x+10} \qquad h(x)=\sqrt{x-7} \qquad
  k(x)=\sqrt{3x+9}
\]
\reponse{$D_{f}=\R\setminus\{6\}$ ;\; $D_{g}=\R\setminus\{-2\}$ ;\;
$D_{h}=\left[7\,;+\infty\right[$ ;\; $D_{k}=\left[-3\,;+\infty\right[$.}
\end{methodebox}

\afaire{exercices 5 à 8}

\partiecours{Tableau de valeurs, courbe et lecture graphique}

\begin{definitionbox}[Définition 3 --- Courbe représentative]
La \textbf{courbe représentative} $\Cf$ de $f$ est formée des points $M\bigl(x\,;f(x)\bigr)$.
Autrement dit, le point $M(x\,;y)$ est sur $\Cf$ si et seulement si $y=f(x)$.
\end{definitionbox}

\begin{methodebox}[Méthode 3 --- Tracer une courbe à partir d'un tableau de valeurs]
\etapeexemples
Tracer la courbe de la fonction $f$ définie sur $\intff{-2}{4}$ par $f(x)=x^{2}-2x-3$.
\begin{correction}
\cas{Résolution}\par\nopagebreak
\begin{minipage}[c]{0.58\linewidth}\raggedright
\explique{On calcule des images :}{un tableau de valeurs}
\begin{center}\small
\renewcommand{\arraystretch}{1.25}
\begin{tabular}{|c|c|c|c|c|c|c|c|}
\hline
$x$ & $-2$ & $-1$ & $0$ & $1$ & $2$ & $3$ & $4$\\
\hline
$f(x)$ & $5$ & $0$ & $-3$ & $-4$ & $-3$ & $0$ & $5$\\
\hline
\end{tabular}
\end{center}
\explique{On place les points $\bigl(x\,;f(x)\bigr)$, puis on les relie par une courbe
régulière, sans règle.}{par exemple $(1\,;-4)$, le point le plus bas}
\end{minipage}\hfill
\begin{minipage}[c]{0.38\linewidth}\centering
\begin{tikzpicture}[x=0.38cm,y=0.38cm]
  \repere{-3}{-5}{5}{5}
  \gradx{-2,-1,1,2,3,4} \grady{-4,-3,-2,-1,1,2,3,4}
  \draw[coulCourbe,line width=1.1pt,domain=-2:4,samples=60] plot (\x,{\x*\x-2*\x-3});
  \foreach \p in {(-2,5),(-1,0),(0,-3),(1,-4),(2,-3),(3,0),(4,5)}
    \path[coulCourbe] \p node[croix]{};
  \node[coulCourbe,font=\footnotesize] at (4.3,3) {$\Cf$};
\end{tikzpicture}
\end{minipage}
\end{correction}
\end{methodebox}

\begin{methodebox}[Méthode 4 --- Lire une image, lire des antécédents]
\etapeexemples
Sur la courbe de la méthode 3 :
\begin{questions}[label=\alph*)]
  \item lire l'image de $1$ ;
  \item lire les antécédents de $0$.
\end{questions}
\begin{correction}
\begin{minipage}[c]{0.58\linewidth}\raggedright
\cas{a)}
\explique{$f(1)=-4$.}{de $1$ sur l'axe des abscisses, on va verticalement jusqu'à la
courbe, puis on lit l'ordonnée (en vert)}
\cas{b)}
\explique{Les antécédents de $0$ sont $-1$ et $3$.}{on trace l'horizontale $y=0$ et on
lit les abscisses de \textbf{tous} les points de la courbe qui sont dessus (en rouge)}
\end{minipage}\hfill
\begin{minipage}[c]{0.38\linewidth}\centering
\begin{tikzpicture}[x=0.38cm,y=0.38cm]
  \repere{-3}{-5}{5}{5}
  \gradx{-2,-1,1,2,3,4} \grady{-4,-3,-2,-1,1,2,3,4}
  \draw[coulCourbe,line width=1.1pt,domain=-2:4,samples=60] plot (\x,{\x*\x-2*\x-3});
  \draw[coulEx,dashed,thick] (1,0) -- (1,-4) -- (0,-4);
  \path[coulEx] (1,-4) node[croix]{};
  \path[coulExo] (-1,0) node[croix]{} (3,0) node[croix]{};
  \node[coulCourbe,font=\footnotesize] at (4.3,3) {$\Cf$};
\end{tikzpicture}
\end{minipage}
\end{correction}

\etape{À vous de jouer}
Sur la même courbe, lire :
\begin{questions}
  \item l'image de $-2$, puis celle de $2$ ;
  \item les antécédents de $-3$ ;
  \item les antécédents de $-4$, puis ceux de $-5$.
\end{questions}
\reponse{1\textdegree) $f(-2)=5$ et $f(2)=-3$ ;\; 2\textdegree) $0$ et $2$ ;\;
3\textdegree) un seul antécédent de $-4$ : $1$ ; aucun pour $-5$, la courbe ne descend
pas jusque-là.}
\end{methodebox}

\afaire{exercices 9 à 13}

\partiecours{Tableau de variations d'une fonction}

\begin{definitionbox}[Définition 4 --- Fonction croissante, décroissante]
Sur un intervalle $I$ :
\begin{itemize}
  \item $f$ est \textbf{croissante} si, quand $x$ augmente, $f(x)$ augmente : la courbe
        monte ; l'ordre est conservé ;
  \item $f$ est \textbf{décroissante} si, quand $x$ augmente, $f(x)$ diminue : la courbe
        descend ; l'ordre est renversé.
\end{itemize}
\begin{center}
\begin{minipage}{0.46\linewidth}\centering
{\small\textbf{$f$ croissante}}\par\smallskip
\begin{tikzpicture}[x=0.6cm,y=0.6cm]
  \draw[grille,step=1] (0,0) grid (4,4);
  \draw[-{Stealth[length=2mm]},thick] (0,0) -- (4.5,0) node[below left=-1pt] {$x$};
  \draw[-{Stealth[length=2mm]},thick] (0,0) -- (0,4.5) node[below left=-1pt] {$y$};
  \draw[coulCourbe,line width=1.1pt,domain=0.2:4,samples=40] plot (\x,{0.4+0.22*\x*\x});
  \draw[coulEx,dashed] (1,0) -- (1,0.62) -- (0,0.62);
  \draw[coulEx,dashed] (3,0) -- (3,2.38) -- (0,2.38);
  \path[coulEx] (1,0.62) node[croix]{} (3,2.38) node[croix]{};
  \node[below,font=\scriptsize] at (1,0) {$x_{1}$};
  \node[below,font=\scriptsize] at (3,0) {$x_{2}$};
  \node[left,font=\scriptsize] at (0,0.62) {$f(x_{1})$};
  \node[left,font=\scriptsize] at (0,2.38) {$f(x_{2})$};
\end{tikzpicture}\par
{\footnotesize $x_{1}<x_{2}$ donne $f(x_{1})<f(x_{2})$}
\end{minipage}\hfill
\begin{minipage}{0.46\linewidth}\centering
{\small\textbf{$f$ décroissante}}\par\smallskip
\begin{tikzpicture}[x=0.6cm,y=0.6cm]
  \draw[grille,step=1] (0,0) grid (4,4);
  \draw[-{Stealth[length=2mm]},thick] (0,0) -- (4.5,0) node[below left=-1pt] {$x$};
  \draw[-{Stealth[length=2mm]},thick] (0,0) -- (0,4.5) node[below left=-1pt] {$y$};
  \draw[coulCourbe,line width=1.1pt,domain=0.2:4,samples=40] plot (\x,{3.9-0.22*\x*\x});
  \draw[coulEx,dashed] (1,0) -- (1,3.68) -- (0,3.68);
  \draw[coulEx,dashed] (3,0) -- (3,1.92) -- (0,1.92);
  \path[coulEx] (1,3.68) node[croix]{} (3,1.92) node[croix]{};
  \node[below,font=\scriptsize] at (1,0) {$x_{1}$};
  \node[below,font=\scriptsize] at (3,0) {$x_{2}$};
  \node[left,font=\scriptsize] at (0,3.68) {$f(x_{1})$};
  \node[left,font=\scriptsize] at (0,1.92) {$f(x_{2})$};
\end{tikzpicture}\par
{\footnotesize $x_{1}<x_{2}$ donne $f(x_{1})>f(x_{2})$}
\end{minipage}
\end{center}
\end{definitionbox}

\begin{definitionbox}[Définition 5 --- Tableau de variations, extremums]
Le \textbf{tableau de variations} résume le sens de variation : une flèche qui monte
quand $f$ est croissante, une flèche qui descend quand elle est décroissante, et les
valeurs de $f$ au bout de chaque flèche.

\smallskip
Sur un intervalle, le \textbf{maximum} de $f$ est la plus grande valeur de $f(x)$, le
\textbf{minimum} la plus petite. On dit toujours \textbf{en quel $x$} il est atteint.
\end{definitionbox}

\begin{methodebox}[Méthode 5 --- Dresser un tableau de variations à partir d'une courbe]
\etapeexemples
La fonction $f$ est définie sur $\intff{-4}{4}$ par sa courbe ci-dessous. Dresser son
tableau de variations, puis donner son maximum et son minimum.
\begin{center}
\begin{tikzpicture}[x=0.42cm,y=0.42cm]
  \repere{-5}{-4}{5}{4}
  \gradx{-4,-3,-2,-1,1,2,3,4} \grady{-3,-2,-1,1,2,3}
  \draw[coulCourbe,line width=1.1pt]
    (-4,-3) .. controls (-3.4,-0.6) and (-2.6,3) .. (-2,3)
            .. controls (-1,3) and (1,-2) .. (2,-2)
            .. controls (2.8,-2) and (3.6,0.4) .. (4,2);
  \foreach \p in {(-4,-3),(-2,3),(2,-2),(4,2)} \path[coulCourbe] \p node[croix]{};
  \node[coulCourbe,font=\footnotesize] at (4.4,3) {$\Cf$};
\end{tikzpicture}
\end{center}
\begin{correction}
\cas{Résolution}\par\nopagebreak
\explique{La courbe monte de $x=-4$ à $x=-2$, descend jusqu'à $x=2$, puis remonte
jusqu'à $x=4$.}{on repère où la courbe change de sens}
\explique{$f(-4)=-3$ ; \ $f(-2)=3$ ; \ $f(2)=-2$ ; \ $f(4)=2$.}{on lit les valeurs aux
bornes et aux changements de sens}
\begin{center}
\begin{tikzpicture}
  \tkzTabInit[lgt=1.6,espcl=2.4]{$x$/1,$f$/2}{$-4$,$-2$,$2$,$4$}
  \tkzTabVar{-/$-3$,+/$3$,-/$-2$,+/$2$}
\end{tikzpicture}
\end{center}
\explique{Maximum : $3$, atteint en $x=-2$. \\ Minimum : $-3$, atteint en $x=-4$.}{on
compare toutes les valeurs du bas des flèches ; le minimum peut être à une borne}
\end{correction}

\etape{À vous de jouer}
$g$ est définie sur $\intff{-5}{3}$ ; elle est décroissante sur $\intff{-5}{-1}$ et
croissante sur $\intff{-1}{3}$, avec $g(-5)=4$, $g(-1)=-2$ et $g(3)=3$.
\begin{questions}
  \item Dresser le tableau de variations de $g$.
  \item Donner le maximum et le minimum de $g$.
\end{questions}
\reponse{1\textdegree) flèche qui descend de $4$ à $-2$, puis qui monte jusqu'à $3$ ;\;
2\textdegree) maximum $4$ en $x=-5$ ; minimum $-2$ en $x=-1$.}
\end{methodebox}

\afaire{exercices 14 à 19}

\partiecours{Fonctions affines --- équation réduite $y=mx+p$}

\begin{definitionbox}[Définition 6 --- Fonction affine]
Une fonction \textbf{affine} s'écrit $f(x)=mx+p$, où $m$ et $p$ sont deux nombres. Elle
est définie sur $\R$, et sa courbe est la droite d'\textbf{équation réduite} $y=mx+p$ :
\begin{itemize}
  \item $m$ est le \textbf{coefficient directeur}, ou \textbf{pente} ;
  \item $p$ est l'\textbf{ordonnée à l'origine} : $f(0)=p$, la droite coupe l'axe des
        ordonnées en $(0\,;p)$.
\end{itemize}
Si $p=0$, $f(x)=mx$ est \textbf{linéaire} ; si $m=0$, $f(x)=p$ est \textbf{constante}.
\end{definitionbox}

\begin{proprietebox}[Propriété 2 --- Sens de variation, pente]
Si $m>0$, $f$ est croissante ; si $m<0$, décroissante ; si $m=0$, constante. Seul le
signe de $m$ compte.

\smallskip
Si $A(x_{A}\,;y_{A})$ et $B(x_{B}\,;y_{B})$ sont deux points de la droite, avec
$x_{A}\ne x_{B}$ :
\[
  m=\frac{y_{B}-y_{A}}{x_{B}-x_{A}}
\]
\end{proprietebox}

\begin{methodebox}[Méthode 6 --- Déterminer une équation réduite]
\etapeexemples
\begin{questions}[label=\alph*)]
  \item Déterminer l'équation réduite de la droite $(AB)$, avec $A(-1\,;4)$ et
        $B(3\,;-4)$.
  \item Déterminer l'équation réduite de la droite de pente $3$ qui passe par
        $C(2\,;1)$.
\end{questions}
\begin{correction}
\cas{a)}
\begin{calculs}
m &= \frac{y_{B}-y_{A}}{x_{B}-x_{A}} &\qquad& \rem{formule de la pente : $B$ en premier, en haut et en bas}\\
m &= \frac{-4-4}{3-(-1)} &&\\
m &= -2 &\qquad& \rem{$\frac{-8}{4}=-2$}
\end{calculs}
\begin{calculs}
y_{A} &= m\,x_{A}+p &&\\
4 &= -2\times(-1)+p &\qquad& \rem{$A$ est sur la droite : ses coordonnées vérifient $y=mx+p$}\\
4 &= 2+p &&\\
p &= 2 &&
\end{calculs}
\explique{$(AB)$ a pour équation réduite $y=-2x+2$.}{contrôle avec $B$ : $-2\times 3+2=-4$}
\cas{b)}
\begin{calculs}
1 &= 3\times 2+p &&\\
p &= -5 &\qquad& \rem{$m=3$ ; $C$ est sur la droite}
\end{calculs}
\explique{L'équation réduite est $y=3x-5$.}{}
\end{correction}

\etape{À vous de jouer}
\begin{questions}
  \item Déterminer l'équation réduite de $(AB)$, avec $A(-2\,;-5)$ et $B(2\,;3)$.
  \item Déterminer l'équation réduite de la droite de pente $-\dfrac{1}{2}$ qui passe par
        $D(4\,;-1)$.
\end{questions}
\reponse{1\textdegree) $m=\frac{8}{4}=2$, puis $p=-1$ : $y=2x-1$ ;\;
2\textdegree) $-1=-2+p$, donc $p=1$ : $y=-\frac{1}{2}x+1$.}
\end{methodebox}

\afaire{exercices 20 à 23}

\partiecours{Tracer la droite d'équation $y=mx+p$}

Pour tracer une droite, deux points suffisent. Le plus rapide : on place $(0\,;p)$ sur
l'axe des ordonnées, puis on avance de $1$ et on monte de $m$ (on descend si $m<0$).

\begin{methodebox}[Méthode 7 --- Tracer une droite, lire son équation]
\etapeexemples
\begin{questions}[label=\alph*)]
  \item Tracer les droites $d_{1}$ : $y=3x-2$ et $d_{2}$ : $y=-\dfrac{2}{3}x+3$.
  \item Une droite coupe l'axe des ordonnées en $(0\,;2)$ ; quand on avance de $2$, elle
        descend de $3$. Donner son équation réduite.
\end{questions}
\begin{correction}
\begin{minipage}[c]{0.58\linewidth}\raggedright
\cas{a)}
\explique{$d_{1}$ : on part de $(0\,;-2)$ ; on avance de $1$, on monte de $3$ : point
$(1\,;1)$.}{$p=-2$ et $m=3$}
\explique{$d_{2}$ : on part de $(0\,;3)$ ; on avance de $3$, on descend de $2$ : point
$(3\,;1)$.}{$m=\frac{-2}{3}$ : avancer de $3$ fait tomber sur un nœud du quadrillage}
\cas{b)}
\explique{$p=2$ et $m=\dfrac{-3}{2}$ : $y=-\dfrac{3}{2}x+2$.}{$m=\frac{\text{déplacement
vertical}}{\text{déplacement horizontal}}$}
\end{minipage}\hfill
\begin{minipage}[c]{0.38\linewidth}\centering
\begin{tikzpicture}[x=0.38cm,y=0.38cm]
  \repere{-5}{-5}{5}{5}
  \gradx{-4,-3,-2,-1,1,2,3,4} \grady{-4,-3,-2,-1,1,2,3,4}
  \draw[coulCourbe,line width=1.1pt,domain=-1:7/3] plot (\x,{3*\x-2});
  \draw[coulCourbeB,line width=1.1pt,domain=-3:5] plot (\x,{-2/3*\x+3});
  \draw[coulCourbe,dashed] (0,-2) -- (1,-2) -- (1,1);
  \draw[coulCourbeB,dashed] (0,3) -- (3,3) -- (3,1);
  \path[coulCourbe] (0,-2) node[croix]{} (1,1) node[croix]{};
  \path[coulCourbeB] (0,3) node[croix]{} (3,1) node[croix]{};
  \node[coulCourbe,font=\footnotesize,right] at (2.3,4.6) {$d_{1}$};
  \node[coulCourbeB,font=\footnotesize] at (4.2,0.6) {$d_{2}$};
\end{tikzpicture}
\end{minipage}
\end{correction}

\etape{À vous de jouer}
\begin{questions}
  \item Tracer les droites d'équations $y=2x+1$ et $y=-x+3$.
  \item Une droite coupe l'axe des ordonnées en $(0\,;-2)$ ; quand on avance de $4$, elle
        monte de $3$. Donner son équation réduite.
\end{questions}
\reponse{1\textdegree) points $(0\,;1)$ et $(1\,;3)$ ; points $(0\,;3)$ et $(3\,;0)$ ;\;
2\textdegree) $y=\frac{3}{4}x-2$.}
\end{methodebox}

\afaire{exercices 24 à 27}

\partiecours{Équation cartésienne d'une droite}

\begin{definitionbox}[Définition 7 --- Équation cartésienne]
Toute droite a une \textbf{équation cartésienne} $ax+by+c=0$, avec $a$ et $b$ non tous
les deux nuls : le point $M(x\,;y)$ est sur la droite si et seulement si ses
coordonnées vérifient l'égalité.
\end{definitionbox}

\begin{proprietebox}[Propriété 3 --- Vecteur directeur]
La droite d'équation $ax+by+c=0$ a pour \textbf{vecteur directeur} $\vec{u}\,(-b\,;a)$.
\end{proprietebox}

\begin{remarquebox}
Une droite a une infinité d'équations cartésiennes : $3x-y+2=0$ et $6x-2y+4=0$ (tout
multiplié par $2$) décrivent la même droite. Son équation réduite, elle, est unique,
mais une droite verticale n'en a pas.
\end{remarquebox}

\begin{methodebox}[Méthode 8 --- Trouver une équation cartésienne]
\etapeexemples
\begin{questions}[label=\alph*)]
  \item Déterminer une équation cartésienne de la droite $d$ qui passe par $A(1\,;4)$
        et de vecteur directeur $\vec{u}\,(3\,;2)$.
  \item Déterminer une équation cartésienne de la droite $(BC)$, avec $B(-2\,;1)$ et
        $C(3\,;-3)$.
\end{questions}
\begin{correction}
\cas{a)}
\explique{$\vec{u}\,(3\,;2)=\vec{u}\,(-b\,;a)$ : $b=-3$ et $a=2$. \\ $d$ : $2x-3y+c=0$.}{%
on identifie les coordonnées de $\vec{u}$ à $(-b\,;a)$}
\begin{calculs}
2\times 1-3\times 4+c &= 0 &&\\
-10+c &= 0 &\qquad& \rem{$A$ est sur $d$ : on remplace $x$ et $y$}\\
c &= 10 &&
\end{calculs}
\explique{$d$ : $2x-3y+10=0$.}{}
\cas{b)}
\begin{calculs}
\vec{BC}\,\bigl(x_{C}-x_{B}\,;\,y_{C}-y_{B}\bigr) &= \vec{BC}\,\bigl(3-(-2)\,;\,-3-1\bigr) &\qquad& \rem{vecteur directeur : « arrivée moins départ »}\\
\vec{BC}\,\bigl(x_{C}-x_{B}\,;\,y_{C}-y_{B}\bigr) &= \vec{BC}\,(5\,;-4) &&
\end{calculs}
\explique{$b=-5$ et $a=-4$ : $-4x-5y+c=0$.}{$(5\,;-4)=(-b\,;a)$}
\begin{calculs}
-4\times(-2)-5\times 1+c &= 0 &&\\
3+c &= 0 &\qquad& \rem{$B$ est sur la droite}\\
c &= -3 &&
\end{calculs}
\explique{$(BC)$ : $-4x-5y-3=0$, ou encore $4x+5y+3=0$.}{contrôle avec $C$ :
$12-15+3=0$}
\end{correction}

\etape{À vous de jouer}
\begin{questions}
  \item Déterminer une équation cartésienne de la droite qui passe par $A(2\,;-1)$ et de
        vecteur directeur $\vec{u}\,(1\,;4)$.
  \item Déterminer une équation cartésienne de $(EF)$, avec $E(1\,;2)$ et $F(3\,;8)$.
\end{questions}
\reponse{1\textdegree) $4x-y+c=0$ avec $c=-9$ : $4x-y-9=0$ ;\;
2\textdegree) $\vec{EF}\,(2\,;6)$ : $6x-2y+c=0$ avec $c=-2$, soit $6x-2y-2=0$ ou
$3x-y-1=0$.}
\end{methodebox}

\afaire{exercices 28 à 30}

\partiecours{Tracer une droite d'équation $ax+by+c=0$}

Pour tracer une droite d'équation cartésienne, on cherche deux de ses points : on choisit
une valeur de $x$ (souvent $x=0$), on calcule $y$, puis on recommence. Si $b=0$, la
droite est \textbf{verticale} ($x=$ nombre) ; si $a=0$, elle est \textbf{horizontale}
($y=$ nombre).

\begin{methodebox}[Méthode 9 --- Tracer une droite d'équation cartésienne]
\etapeexemples
Tracer dans un même repère les droites $d_{1}$ : $x+2y-4=0$, $d_{2}$ : $3x-y-2=0$ et
$d_{3}$ : $2x+6=0$.
\begin{correction}
\begin{minipage}[c]{0.58\linewidth}\raggedright
\cas{Résolution}\par\nopagebreak
\explique{$d_{1}$ : pour $x=0$ on a $2y-4=0$, donc $y=2$ ; pour $y=0$, $x=4$. Points
$(0\,;2)$ et $(4\,;0)$.}{on remplace $x$ (ou $y$) par $0$ : les points sur les axes}
\explique{$d_{2}$ : pour $x=0$, $y=-2$ ; pour $x=2$, $6-y-2=0$, donc $y=4$. Points
$(0\,;-2)$ et $(2\,;4)$.}{pour $y=0$, $x=\frac{2}{3}$ tombe mal : on choisit une autre
valeur de $x$}
\explique{$d_{3}$ : $2x+6=0$ donne $x=-3$.}{pas de $y$ : droite verticale}
\end{minipage}\hfill
\begin{minipage}[c]{0.38\linewidth}\centering
\begin{tikzpicture}[x=0.38cm,y=0.38cm]
  \repere{-5}{-5}{5}{5}
  \gradx{-4,-3,-2,-1,1,2,3,4} \grady{-4,-3,-2,-1,1,2,3,4}
  \draw[coulCourbe,line width=1.1pt,domain=-2:5] plot (\x,{2-\x/2});
  \draw[coulCourbeB,line width=1.1pt,domain=-1:7/3] plot (\x,{3*\x-2});
  \draw[coulCourbeC,line width=1.1pt] (-3,-5) -- (-3,5);
  \path[coulCourbe] (0,2) node[croix]{} (4,0) node[croix]{};
  \path[coulCourbeB] (0,-2) node[croix]{} (2,4) node[croix]{};
  \node[coulCourbe,font=\footnotesize] at (4.5,0.8) {$d_{1}$};
  \node[coulCourbeB,font=\footnotesize] at (1.4,4.6) {$d_{2}$};
  \node[coulCourbeC,font=\footnotesize] at (-3.9,4.4) {$d_{3}$};
\end{tikzpicture}
\end{minipage}
\end{correction}

\etape{À vous de jouer}
Tracer dans un même repère les droites d'équations $x-y+1=0$ ; $2x+y-4=0$ et $y-3=0$.
\reponse{points $(0\,;1)$ et $(-1\,;0)$ ; points $(0\,;4)$ et $(2\,;0)$ ; droite
horizontale $y=3$.}
\end{methodebox}

\afaire{exercices 31 à 33}

\partiecours{Passer d'une forme à l'autre}

\begin{proprietebox}[Propriété 4 --- Équation réduite et équation cartésienne]
\begin{itemize}
  \item De la réduite à la cartésienne : on met tout dans un membre,
        $y=mx+p$ devient $mx-y+p=0$.
  \item De la cartésienne à la réduite : on isole $y$, comme dans une équation ; c'est
        possible seulement si $b\ne 0$.
\end{itemize}
\end{proprietebox}

\begin{methodebox}[Méthode 10 --- Passer d'une forme à l'autre]
\etapeexemples
\begin{questions}[label=\alph*)]
  \item Donner une équation cartésienne de la droite d'équation $y=\dfrac{2}{3}x-1$.
  \item Donner l'équation réduite de la droite d'équation $4x+2y-6=0$, sa pente et son
        ordonnée à l'origine.
  \item La droite d'équation $3x+9=0$ a-t-elle une équation réduite ?
\end{questions}
\begin{correction}
\cas{a)}
\begin{calculs}
y &= \dfrac{2}{3}x-1 &&\\
\dfrac{2}{3}x-y-1 &= 0 &\qquad& \rem{tout dans le membre de gauche}\\
2x-3y-3 &= 0 &\qquad& \rem{on multiplie par $3$ pour ôter la fraction}
\end{calculs}
\cas{b)}
\begin{calculs}
4x+2y-6 &= 0 &&\\
2y &= -4x+6 &\qquad& \rem{on isole $2y$}\\
y &= -2x+3 &\qquad& \rem{on divise \textbf{les deux termes} par $2$}
\end{calculs}
\explique{Pente $m=-2$, ordonnée à l'origine $p=3$.}{}
\cas{c)}
\explique{$3x+9=0$ donne $x=-3$ : droite verticale, pas d'équation réduite.}{$b=0$ : on
ne peut pas isoler $y$ ; ce n'est pas la courbe d'une fonction}
\end{correction}

\etape{À vous de jouer}
\begin{questions}
  \item Donner une équation cartésienne de la droite $y=-3x+4$.
  \item Donner l'équation réduite de la droite $2x-5y+10=0$, puis sa pente.
  \item Les droites $2x-y+3=0$ et $y=2x+3$ sont-elles confondues ?
\end{questions}
\reponse{1\textdegree) $3x+y-4=0$ ;\; 2\textdegree) $y=\frac{2}{5}x+2$, pente
$\frac{2}{5}$ ;\; 3\textdegree) oui : $2x-y+3=0$ donne $y=2x+3$.}
\end{methodebox}

\afaire{exercices 34 à 37}

\partiecours{Fonction carré}

\begin{definitionbox}[Définition 8 --- Fonction carré]
La \textbf{fonction carré} est définie sur $\R$ par $f(x)=x^{2}$. Sa courbe est une
\textbf{parabole} de sommet l'origine $O$.
\end{definitionbox}

\begin{center}
\begin{minipage}{0.40\linewidth}\centering
\begin{tikzpicture}[x=0.42cm,y=0.42cm]
  \repere{-3}{-1}{3}{9}
  \gradx{-2,-1,1,2} \grady{1,2,3,4,5,6,7,8}
  \draw[coulCourbe,line width=1.1pt,domain=-3:3,samples=60] plot (\x,{\x*\x});
  \foreach \p in {(-3,9),(-2,4),(-1,1),(0,0),(1,1),(2,4),(3,9)} \path[coulCourbe] \p node[croix]{};
\end{tikzpicture}
\end{minipage}\hfill
\begin{minipage}{0.55\linewidth}\centering
\begin{tikzpicture}[scale=0.85]
  \tkzTabInit[lgt=1.6,espcl=1.8]{$x$/1,$x^{2}$/2}{$-\infty$,$0$,$+\infty$}
  \tkzTabVar{+/$+\infty$,-/$0$,+/$+\infty$}
\end{tikzpicture}

\begin{proprietebox}[Propriété 5 --- Carré]
Décroissante sur $\left]-\infty\,;0\right]$, croissante sur $\left[0\,;+\infty\right[$.
Un carré est toujours positif ou nul, et deux nombres opposés ont le même carré :
$(-3)^{2}=3^{2}=9$.
\end{proprietebox}
\end{minipage}
\end{center}

\begin{proprietebox}[Propriété 6 --- Équation $x^{2}=k$]
Si $k>0$ : deux solutions, $-\sqrt{k}$ et $\sqrt{k}$. Si $k=0$ : une solution, $0$. Si
$k<0$ : aucune solution.
\end{proprietebox}

\begin{methodebox}[Méthode 11 --- Résoudre une équation en isolant $x^{2}$]
\etapeexemples
Résoudre dans $\R$ les équations suivantes.
\begin{questions}[label=\alph*)]
  \item $x^{2}=16$
  \item $5x^{2}+1=2x^{2}+16$
  \item $3x^{2}+8=5$
\end{questions}
\begin{correction}
\cas{a)}
\explique{$S=\ens{-4\,;4}$.}{$16>0$ : deux solutions ; on n'oublie pas $-4$, car
$(-4)^{2}=16$}
\cas{b)}
\begin{calculs}
5x^{2}+1 &= 2x^{2}+16 &&\\
5x^{2}-2x^{2} &= 16-1 &\qquad& \rem{on isole $x^{2}$ comme on isole $x$}\\
3x^{2} &= 15 &&\\
x^{2} &= 5 &&
\end{calculs}
\explique{$S=\ens{-\sqrt{5}\,;\sqrt{5}}$.}{$5>0$ : deux solutions}
\cas{c)}
\begin{calculs}
3x^{2}+8 &= 5 &&\\
3x^{2} &= -3 &&\\
x^{2} &= -1 &&
\end{calculs}
\explique{$S=\varnothing$.}{un carré n'est jamais négatif}
\end{correction}

\etape{À vous de jouer}
Résoudre dans $\R$ :
\[
  x^{2}=64 \qquad 4x^{2}-5=2x^{2}+9 \qquad 3x^{2}+1=13 \qquad 2x^{2}+9=1
\]
\reponse{$S=\{-8\,;8\}$ ;\; $x^{2}=7$, $S=\left\{-\sqrt{7}\,;\sqrt{7}\right\}$ ;\;
$x^{2}=4$, $S=\{-2\,;2\}$ ;\; $x^{2}=-4$, $S=\varnothing$.}
\end{methodebox}

\afaire{exercices 38 à 40}

\partiecours{Fonction cube}

\begin{definitionbox}[Définition 9 --- Fonction cube]
La \textbf{fonction cube} est définie sur $\R$ par $g(x)=x^{3}$. Sa courbe a pour centre
de symétrie l'origine $O$.
\end{definitionbox}

\begin{center}
\begin{minipage}{0.40\linewidth}\centering
\begin{tikzpicture}[x=0.42cm,y=0.42cm]
  \repere{-3}{-5}{3}{5}
  \gradx{-2,-1,1,2} \grady{-4,-3,-2,-1,1,2,3,4}
  \begin{scope}
    \clip (-3,-5) rectangle (3,5);
    \draw[coulCourbe,line width=1.1pt,domain=-1.72:1.72,samples=60] plot (\x,{\x*\x*\x});
  \end{scope}
  \foreach \p in {(-1,-1),(0,0),(1,1)} \path[coulCourbe] \p node[croix]{};
\end{tikzpicture}
\end{minipage}\hfill
\begin{minipage}{0.55\linewidth}\centering
\begin{tikzpicture}[scale=0.85]
  \tkzTabInit[lgt=1.6,espcl=3.4]{$x$/1,$x^{3}$/2}{$-\infty$,$+\infty$}
  \tkzTabVar{-/$-\infty$,+/$+\infty$}
\end{tikzpicture}

\smallskip
{\footnotesize
\renewcommand{\arraystretch}{1.2}
\begin{tabular}{|c|c|c|c|c|c|c|c|}
\hline
$x$ & $-3$ & $-2$ & $-1$ & $0$ & $1$ & $2$ & $3$\\
\hline
$x^{3}$ & $-27$ & $-8$ & $-1$ & $0$ & $1$ & $8$ & $27$\\
\hline
\end{tabular}}
\end{minipage}
\end{center}

\begin{proprietebox}[Propriété 7 --- Cube]
La fonction cube est croissante sur $\R$. Un cube a le signe du nombre :
$(-2)^{3}=-8$ et $2^{3}=8$.

\smallskip
Pour tout nombre $k$, l'équation $x^{3}=k$ a \textbf{une seule} solution, la
\textbf{racine cubique} de $k$, notée $\sqrt[3]{k}$ (touche de la calculatrice).
\end{proprietebox}

\begin{methodebox}[Méthode 12 --- Résoudre une équation en isolant $x^{3}$]
\etapeexemples
Résoudre dans $\R$ les équations suivantes.
\begin{questions}[label=\alph*)]
  \item $x^{3}=-27$
  \item $4x^{3}-5=x^{3}+19$
  \item $2x^{3}+9=5x^{3}-6$
\end{questions}
\begin{correction}
\cas{a)}
\explique{$S=\ens{-3}$.}{$(-3)^{3}=-27$ ; une seule solution, même pour un nombre
négatif}
\cas{b)}
\begin{calculs}
4x^{3}-5 &= x^{3}+19 &&\\
3x^{3} &= 24 &\qquad& \rem{on isole $x^{3}$}\\
x^{3} &= 8 &&
\end{calculs}
\explique{$S=\ens{2}$.}{$2^{3}=8$}
\cas{c)}
\begin{calculs}
2x^{3}+9 &= 5x^{3}-6 &&\\
-3x^{3} &= -15 &&\\
x^{3} &= 5 &&
\end{calculs}
\explique{$S=\ens{\sqrt[3]{5}}$.}{$\sqrt[3]{5}\approx 1{,}71$ à la calculatrice}
\end{correction}

\etape{À vous de jouer}
Résoudre dans $\R$ :
\[
  x^{3}=125 \qquad 3x^{3}=-3 \qquad 5x^{3}+2=2x^{3}-22 \qquad 7x^{3}-4=x^{3}+14
\]
\reponse{$S=\{5\}$ ;\; $x^{3}=-1$, $S=\{-1\}$ ;\; $x^{3}=-8$, $S=\{-2\}$ ;\;
$x^{3}=3$, $S=\left\{\sqrt[3]{3}\right\}$.}
\end{methodebox}

\afaire{exercices 41 à 43}

\partiecours{Fonction inverse}

\begin{definitionbox}[Définition 10 --- Fonction inverse]
La \textbf{fonction inverse} est définie pour tout $x\ne 0$ par $f(x)=\dfrac{1}{x}$ :
$D_{f}=\R\setminus\{0\}$. Sa courbe est une \textbf{hyperbole}, de centre de symétrie $O$.
\end{definitionbox}

\begin{center}
\begin{minipage}{0.40\linewidth}\centering
\begin{tikzpicture}[x=0.42cm,y=0.42cm]
  \repere{-4}{-4}{4}{4}
  \gradx{-3,-2,-1,1,2,3} \grady{-3,-2,-1,1,2,3}
  \draw[coulCourbe,line width=1.1pt,domain=0.25:4,samples=60] plot (\x,{1/\x});
  \draw[coulCourbe,line width=1.1pt,domain=-4:-0.25,samples=60] plot (\x,{1/\x});
  \path[coulCourbe] (1,1) node[croix]{} (-1,-1) node[croix]{};
\end{tikzpicture}
\end{minipage}\hfill
\begin{minipage}{0.55\linewidth}\centering
\begin{tikzpicture}[scale=0.85]
  \tkzTabInit[lgt=1.6,espcl=2]{$x$/1,$\frac{1}{x}$/2}{$-\infty$,$0$,$+\infty$}
  \tkzTabVar{+/$0$,-D+/$-\infty$/$+\infty$,-/$0$}
\end{tikzpicture}

\smallskip
\begin{tikzpicture}[scale=0.85]
  \tkzTabInit[lgt=1.6,espcl=2]{$x$/1,$\frac{1}{x}$/1}{$-\infty$,$0$,$+\infty$}
  \tkzTabLine{,-,d,+,}
\end{tikzpicture}
\end{minipage}
\end{center}

\begin{proprietebox}[Propriété 8 --- Inverse]
La fonction inverse est décroissante sur $\left]-\infty\,;0\right[$ et décroissante sur
$\left]0\,;+\infty\right[$ ; $\dfrac{1}{x}$ a le signe de $x$.

\smallskip
Elle n'est pas décroissante sur $\R\setminus\{0\}$ tout entier : $-1<1$ et
$\dfrac{1}{-1}<\dfrac{1}{1}$. On sépare toujours les deux intervalles.
\end{proprietebox}

\begin{methodebox}[Méthode 13 --- Équations et inéquations avec un inverse]
\etapeexemples
Résoudre dans $\R$ :
\begin{questions}[label=\alph*)]
  \item $\dfrac{1}{x}=4$
  \item $\dfrac{6}{2x-4}=3$
  \item $\dfrac{1}{x}\le 3$
\end{questions}
\begin{correction}
\cas{a)}
\begin{calculs}
\frac{1}{x} &= 4 &&\\
1 &= 4x &\qquad& \rem{$x\ne 0$ : on multiplie les deux membres par $x$}\\
x &= \frac{1}{4} &&
\end{calculs}
\explique{$S=\ens{\tfrac{1}{4}}$.}{}
\cas{b)}
\explique{Valeur interdite : $2x-4=0$ pour $x=2$.}{on la cherche d'abord}
\begin{calculs}
\frac{6}{2x-4} &= 3 &&\\
6 &= 3(2x-4) &\qquad& \rem{pour $x\ne 2$, on multiplie par $2x-4$}\\
6 &= 6x-12 &&\\
x &= 3 &&
\end{calculs}
\explique{$S=\ens{3}$.}{$3$ n'est pas la valeur interdite}
\cas{c)}
\begin{calculs}
\frac{1}{x} &\le 3 &&\\
\frac{1}{x}-3 &\le 0 &\qquad& \rem{on ne multiplie pas par $x$ : son signe est inconnu}\\
\frac{1-3x}{x} &\le 0 &\qquad& \rem{un seul quotient (chapitre 1)}
\end{calculs}
\begin{center}
\begin{tikzpicture}
  \tkzTabInit[lgt=2.6,espcl=2.4]{$x$/1,$1-3x$/1,$x$/1,$\frac{1-3x}{x}$/1.4}%
    {$-\infty$,$0$,$\frac{1}{3}$,$+\infty$}
  \tkzTabLine{,+,t,+,z,-,}
  \tkzTabLine{,-,z,+,t,+,}
  \tkzTabLine{,-,d,+,z,-,}
\end{tikzpicture}
\end{center}
\explique{$S=\left]-\infty\,;0\right[\cup\left[\tfrac{1}{3}\,;+\infty\right[$.}{$0$
exclu (interdit) ; $\frac{1}{3}$ compris}
\end{correction}

\etape{À vous de jouer}
Résoudre dans $\R$ :
\[
  \dfrac{1}{x}=-5 \qquad \dfrac{4}{-2x+6}=2 \qquad \dfrac{1}{x}>2
\]
\reponse{$S=\left\{-\frac{1}{5}\right\}$ ;\; valeur interdite $3$, $4=-4x+12$ donne
$S=\{2\}$ ;\; $\frac{1-2x}{x}>0$, $S=\left]0\,;\frac{1}{2}\right[$.}
\end{methodebox}

\afaire{exercices 44 à 48}

\partiecours{Fonction racine carrée}

\begin{definitionbox}[Définition 11 --- Fonction racine carrée]
La \textbf{fonction racine carrée} est définie pour tout $x\ge 0$ par $f(x)=\sqrt{x}$ :
$D_{f}=\left[0\,;+\infty\right[$.
\end{definitionbox}

\begin{center}
\begin{minipage}{0.48\linewidth}\centering
\begin{tikzpicture}[x=0.42cm,y=0.42cm]
  \repere{-1}{-1}{9}{4}
  \gradx{1,2,3,4,5,6,7,8} \grady{1,2,3}
  \draw[coulCourbe,line width=1.1pt,domain=0:9,samples=100] plot (\x,{sqrt(\x)});
  \foreach \p in {(0,0),(1,1),(4,2),(9,3)} \path[coulCourbe] \p node[croix]{};
\end{tikzpicture}
\end{minipage}\hfill
\begin{minipage}{0.48\linewidth}\centering
\begin{tikzpicture}[scale=0.85]
  \tkzTabInit[lgt=1.6,espcl=3.4]{$x$/1,$\sqrt{x}$/2}{$0$,$+\infty$}
  \tkzTabVar{-/$0$,+/$+\infty$}
\end{tikzpicture}

\smallskip
{\footnotesize
\renewcommand{\arraystretch}{1.2}
\begin{tabular}{|c|c|c|c|c|c|}
\hline
$x$ & $0$ & $1$ & $4$ & $9$ & $16$\\
\hline
$\sqrt{x}$ & $0$ & $1$ & $2$ & $3$ & $4$\\
\hline
\end{tabular}}
\end{minipage}
\end{center}

\begin{proprietebox}[Propriété 9 --- Racine carrée]
La fonction racine carrée est croissante sur $\left[0\,;+\infty\right[$, et toujours
positive ou nulle. Pour des nombres positifs, elle conserve l'ordre :
si $a\le b$, alors $\sqrt{a}\le\sqrt{b}$, et réciproquement.
\end{proprietebox}

\begin{methodebox}[Méthode 14 --- Équations et inéquations avec une racine carrée]
\etapeexemples
Résoudre :
\begin{questions}[label=\alph*)]
  \item $\sqrt{x}=5$
  \item $\sqrt{x}=-4$
  \item $\sqrt{x}\le 3$
  \item $\sqrt{x}>4$
  \item $\dfrac{18}{\sqrt{x}}=6$
\end{questions}
\begin{correction}
\cas{a)}
\explique{$\sqrt{x}=5$ \\ $x=25$ \\ $S=\ens{25}$.}{$x\ge 0$ toujours ; les deux membres
sont positifs, on élève au carré}
\cas{b)}
\explique{$S=\varnothing$.}{une racine carrée n'est jamais négative (élever au carré
donnerait $16$, qui ne convient pas)}
\cas{c)}
\explique{$\sqrt{x}\le 3$ \\ $0\le x\le 9$ \\ $S=\intff{0}{9}$.}{la racine conserve
l'ordre ; on n'oublie pas $x\ge 0$}
\cas{d)}
\explique{$\sqrt{x}>4$ \\ $x>16$ \\ $S=\left]16\,;+\infty\right[$.}{on élève au carré : l'ordre est conservé}
\cas{e)}
\begin{calculs}
\frac{18}{\sqrt{x}} &= 6 &&\\
18 &= 6\sqrt{x} &\qquad& \rem{$x>0$ : le dénominateur n'est pas nul}\\
\sqrt{x} &= 3 &&\\
x &= 9 &&
\end{calculs}
\explique{$S=\ens{9}$.}{}
\end{correction}

\etape{À vous de jouer}
Résoudre :
\[
  \sqrt{x}=7 \qquad \sqrt{x}\le 5 \qquad \sqrt{x}\ge 2 \qquad \dfrac{5}{x^{2}}=20
\]
\reponse{$S=\{49\}$ ;\; $S=\intff{0}{25}$ ;\; $S=\left[4\,;+\infty\right[$ ;\;
$x^{2}=\frac{1}{4}$, $S=\left\{-\frac{1}{2}\,;\frac{1}{2}\right\}$.}
\end{methodebox}

\afaire{exercices 49 à 52}

\end{document}
